\chapter{Literature Survey / Related Work}
\label{ch:lit}

\section{Existing Methodologies}
Traditional corporate benchmarking has evolved from manual spreadsheet transcription to proprietary financial databases:
\begin{itemize}
    \item \textbf{Manual Spreadsheet Benchmarking}: Analysts manually extract and transcribe quarterly figures into Microsoft Excel models. This foundational framework, formalized by Damodaran \cite{damodaran2012investment}, remains slow, labor-intensive, and prone to human error.
    \item \textbf{Proprietary Financial Terminals}: Bloomberg Terminals and S\&P Capital IQ automate data harvesting and standardization. However, their high annual licensing fees (\$24,000+ per user) and closed database architectures prevent private enterprises from easily uploading double-entry Tally ERP ledger data for direct cross-sectional benchmarking.
\end{itemize}

\section{Machine Learning \& Robust Outlier Detection in Finance}
Corporate financial datasets frequently contain extreme outliers from non-recurring activities, distorting standard mean-based Z-score benchmarking. To resolve this, robust statistics and unsupervised machine learning models are deployed:
\begin{itemize}
    \item \textbf{Tukey's Fences (IQR)}: Non-parametric filter using the Interquartile Range ($IQR = Q_3 - Q_1$) to define bounds as $[Q_1 - k \cdot IQR, \; Q_3 + k \cdot IQR]$ \cite{tukey1977exploratory}. Resilient against extreme skewness as it relies on medians.
    \item \textbf{Modified Z-Scores (MAD)}: Measures statistical dispersion via the Median Absolute Deviation ($MAD = \text{median}(|x_i - \tilde{x}|)$). The Modified Z-Score $M_i$ is formulated as $M_i = 0.6745(x_i - \tilde{x}) / MAD$. Flags outliers exceeding $3.5$ robustly.
    \item \textbf{Isolation Forests (iForest)}: An unsupervised machine learning algorithm \cite{liu2008isolation} that isolates anomalies in multi-dimensional space by randomly partitioning feature values. Anomalies isolate closer to the tree root, resulting in shorter path lengths.
\end{itemize}

\begin{figure}[htbp]
    \centering
    \begin{tikzpicture}[
        node distance=0.8cm and 0.5cm,
        block/.style={rectangle, draw=black, fill=blue!5, text width=4.0cm, text centered, rounded corners, minimum height=1.2cm, font=\scriptsize},
        badblock/.style={rectangle, draw=black, fill=red!5, text width=4.0cm, text centered, rounded corners, minimum height=1.2cm, font=\scriptsize},
        goodblock/.style={rectangle, draw=black, fill=green!5, text width=4.0cm, text centered, rounded corners, minimum height=1.2cm, font=\scriptsize},
        arrow/.style={thick,->,>=stealth}
    ]
        % Left Track (Traditional)
        \node (trad1) [block] {\textbf{Traditional Benchmarking}\\Simple averages, manual spreadsheet entry, static bounds};
        \node (trad2) [badblock, below=of trad1] {\textbf{Outlier Contamination}\\Contaminated data points distort mean and standard deviation};
        \node (trad3) [badblock, below=of trad2] {\textbf{Distorted Baselines}\\Skewed dispersion bounds, subjective credit modeling};

        % Right Track (Advanced Robust/ML)
        \node (ml1) [block, right=1.5cm of trad1] {\textbf{FinMetrics Robust Suite}\\Automated statement parsing, multi-layered filters};
        \node (ml2) [goodblock, below=of ml1] {\textbf{Dynamic Stats \& ML}\\Tukey Fences (IQR), Modified Z-Score (MAD), Isolation Forest};
        \node (ml3) [goodblock, below=of ml2] {\textbf{Clean Peer Baselines}\\Robust, outlier-masked baselines and K-Means archetypes};

        % Connections
        \draw [arrow] (trad1) -- (trad2);
        \draw [arrow] (trad2) -- (trad3);
        \draw [arrow] (ml1) -- (ml2);
        \draw [arrow] (ml2) -- (ml3);
    \end{tikzpicture}
    \caption{Conceptual Comparison of Traditional vs. Robust Machine Learning Benchmarking Pipelines}
    \label{fig:lit_workflow}
\end{figure}

\section{K-Means Clustering \& Corporate Archetypes}
Traditional industry classifications group companies with vastly different financial architectures. Unsupervised clustering offers a robust mathematical alternative:
\begin{itemize}
    \item \textbf{K-Means Centroid Partitioning}: Groups $n$ observations into $K$ clusters by minimizing the Within-Cluster Sum of Squares (WCSS): $\sum\limits_{k=1}^K \sum\limits_{x_i \in C_k} ||x_i - \mu_k||^2$ \cite{macqueen1967some}.
    \item \textbf{Archetype Segmentation}: Clusters are determined via the Elbow Method/Silhouette analysis and mapped to distinct operational categories (e.g., \textit{"Elite Market Leaders"} vs. \textit{"High-Risk / Distressed"}).
\end{itemize}

\begin{figure}[htbp]
    \centering
    % Define soft pastel colors locally inside the figure
    \definecolor{lavenderbg}{RGB}{244, 245, 252}
    \definecolor{peachbg}{RGB}{253, 246, 238}
    \definecolor{greenbg}{RGB}{242, 252, 242}
    \definecolor{boxborder}{RGB}{80, 80, 80}

    \begin{tikzpicture}[
        % Using a serif font and scriptsize to make the figure more compact and professional
        font=\rmfamily\scriptsize,
        >={Stealth[length=5pt, width=3.75pt]},
        % Base styling for all nodes
        box/.style={
            draw=boxborder, rounded corners=4pt, align=center,
            minimum width=3.3cm, minimum height=1.2cm, inner sep=4pt, thick
        },
        lavender/.style={box, fill=lavenderbg},
        peach/.style={box, fill=peachbg},
        green/.style={box, fill=greenbg}
    ]

        % =========================================================================
        % 1. NODES
        % =========================================================================
        
        % Level 1
        \node[lavender] (raw) at (0, 0) {
            \textbf{Raw Financial Ratios}\\
            (30+ calculated metrics)
        };

        % Level 2 (The 3 outlier algorithms)
        \node[peach] (zscore) at (0, -2.4) {
            \textbf{Modified Z-Score}\\
            \textbf{(MAD)}\\
            Measures robust\\
            median dispersion
        };
        \node[peach] (tukey) at (-4.2, -2.4) {
            \textbf{Tukey's Fences (IQR)}\\
            Filters extreme\\
            univariate range skewness
        };
        \node[peach] (iforest) at (4.2, -2.4) {
            \textbf{Isolation Forest}\\
            \textbf{(iForest)}\\
            Screens multi-dimensional\\
            outliers
        };

        % Level 3
        \node[green] (clean) at (0, -4.8) {
            \textbf{Clean Masked Dataset}\\
            (Outliers removed\\from baselines)
        };

        % Level 4
        \node[lavender] (kmeans) at (0, -7.0) {
            \textbf{K-Means Clustering}\\
            (WCSS centroid\\partitioning)
        };

        % Level 5
        \node[green] (archetypes) at (0, -9.2) {
            \textbf{Corporate Archetypes}\\
            (Elite, High-Risk, Efficient)
        };

        % =========================================================================
        % 2. CONNECTIONS (Orthogonal Routing)
        % =========================================================================
        \begin{scope}[thick, draw=black!80]
            
            % --- Top Fork (1 to 3) ---
            % Calculate a point exactly halfway between Level 1 and Level 2
            \coordinate (fork) at ($(raw.south)!0.5!(zscore.north)$);
            
            \draw (raw.south) -- (fork);
            \draw[->, rounded corners=2pt] (fork) -| (tukey.north);
            \draw[->] (fork) -- (zscore.north);
            \draw[->, rounded corners=2pt] (fork) -| (iforest.north);

            % --- Middle Merge (3 to 1) ---
            % Calculate a point exactly halfway between Level 2 and Level 3
            \coordinate (merge) at ($(zscore.south)!0.5!(clean.north)$);
            
            \draw[rounded corners=2pt] (tukey.south) |- (merge);
            \draw[rounded corners=2pt] (iforest.south) |- (merge);
            \draw (zscore.south) -- (merge);
            \draw[->] (merge) -- (clean.north);

            % --- Bottom Vertical Drops ---
            \draw[->] (clean.south) -- (kmeans.north);
            \draw[->] (kmeans.south) -- (archetypes.north);
            
        \end{scope}

    \end{tikzpicture}
    \caption{Multi-layered Robust Outlier Masking and K-Means Clustering Pipeline}
    \label{fig:outlier_cluster_pipeline}
\end{figure}

\section{Research Gaps}
Three critical research gaps remain in current financial intelligence frameworks:
\begin{itemize}
    \setlength{\itemsep}{2pt}
    \setlength{\parskip}{0pt}
    \setlength{\parsep}{0pt}
    \item \textbf{Lack of Nomenclature Autonomy}: Scrapers fail to handle inconsistent reporting names, causing column mismatches or missing values in raw statement parsing.
    \item \textbf{Baseline Outlier Contamination}: Distressed peer data distorts average standard deviations. Standard methods lack automatic statistical masking before baseline generation.
    \item \textbf{Lack of ERP Double-Entry Integration}: Existing benchmarks require manual ratio inputs and lack streamers to process raw, varying-sign debit/credit transaction ledger CSVs.
\end{itemize}
